   \documentclass[12pt]{article}
\usepackage{geometry}
\geometry{margin=1in}
\usepackage{graphicx}
\usepackage{amsmath}
\usepackage{hyperref}
\usepackage{enumitem}
\usepackage{longtable}
\usepackage{array}
\title{Forecasting Economic Impacts of Digital Services on GDP}
\author{Adam Abdellaoui, Lakya Knox, Brandon McLean \\[1ex]
    Master's of Computer Science \\
    Georgia Southern University \\
    \texttt{bm05663@georgiasouthern.edu}
}
\date{\today}

\begin{document}
\maketitle

\begin{abstract}
\noindent
\textbf{Background:} Throughout the world countries use various measures and means to compare productivity, prosperity, and to drive policy. Of the measures used gross domestic product (GDP), is among the most trusted indicators of economic prosperity.  \\
\textbf{Objective:} This study aims to create a prediction model capable of simulating movement in digital economic markets. \\
\textbf{Methods:} Data provided by World Bank, US Census datasets, St. Louis Federal Reserve Economic Data. \\
\textbf{Results:} Predictive analysis, forecasting, digital services, machine-learning. \\
\textbf{Conclusion:} Main takeaway and significance. \\[1ex]
\textbf{Keywords:} machine-learning, predictive analysis, gross domestic product, forecasting
\end{abstract}

% Funding and Conflicts of Interest
\section*{Funding and Conflicts of Interest}
No conflicts of interest, project is unfunded.

\tableofcontents
\newpage

\section{Introduction}
    \subsection*{Background and Motivation}
        From the early 2010s to present day there has been an explosion of wealth generated, primarily in first world economies. This increase in value has largely been made possible by developments in technology and furthermore finance. With an ever more connected world money is circulating faster than ever, with the rise in value of alternative forms of capital such as cryptocurrency the possible gains or losses have never been greater for individuals as well as nations. Predicting best, worst, and average scenarios for GDP would enable policymakers to make more informed decisions to benefit their organization.

        This study aims to investigate the evolving role of digital services in economic growth and the implications for GDP measurement and forecasting. It examines how growth in digital services contributes to overall GDP expansion, taking into account ongoing trends toward service-oriented and platform-based economies. Through empirical analysis and predictive modeling this research aims to bridge gaps in national income measurement and macroeconomic forecasting, the findings are expected to inform policymakers, economists, and data scientist on the basis of improving policy design, decision making, and to enhance forecasting tools to better account for the effects of digitization.
    \subsection*{Problem Definition}
        How does growth in digital services contribute to GDP growth?
        Can ML-based GDP forecasting models be adapted for scenario simulation to evaluate macroeconomic impact of digital services in economies?
    \subsection*{Objectives and Contributions}
        To develop ML models that can simulate events on a macroeconomic scale using digital services and platform data as a focus.

\section{Literature Review}
    \subsection*{Theoretical Background}
        
        Several studies have shown that incorporating additional determinants
such as economic indicators, political events, and oil prices
can improve the predictive accuracy of GDP [1]. More precise estimations in GDP growth, positive or negative, can drive government spending and policy. Several To forecast GDP a wide range of predictive approaches must be employed including: time-series models, linear and non-linear regression, as well as Artificial Neural Networks (ANN). 
        
        Traditional econometric models often struggle when presented with non-linear trends in data, a function that ANN is especially adept at. Richardson, Florenstein Mulder, and Vehbi determined that neural networks, boosted trees, and support vector machine regression algorithms were able to reduce average forecast error rates by 20 - 30\% when compared to auto-regressive (AR) algorithms [2]. Artificial neural networks have been proven to be effective at modeling high-dimensional, non-linear relationships. Vasenka also reports that through a comparative analysis of ML and Deep Machine Learning (DML) achieved results outperforming individual models by 10.2\% [3].
        
        
        This study synthesizes findings from more than 100 peer-reviewed articles (2014 - 2023) examining the application of machine learning (ML) and artificial intelligence (AI) in financial forecasting across four major asset classes: equities, cryptocurrencies, commodities, and foreign exchange markets. Using the Preferred Reporting Items for Systematic Reviews and Meta-Analyses (PRISMA) framework, we have identified Long Short-Term Memory (LSTM) networks, Gated Recurrent Units (GRU), and XGBoost as the most frequently employed predictive models [4]. Google Scholar, Science Direct, and Web of Science served as the source of their datasets. Root Mean Square Error (RMSE) and Mean Absolute Percentage Error (MAPE) emerge as the most reliable standard evaluation metrics. 

        The relationship between digital payment adoption and GDP growth using international macroeconomic data, taking into consideration of behavioral factors such as financial literacy and user attitude, has been studied through traditional econometric panel regression models rather than machine learning techniques. Predictive mean matching, classification and regression trees, as well as Lasso linear regression were implemented through a process known as MICE (Multiple Imputation by Chained Equations), a regression method used to manage missing data by creating multiple complete datasets. Using data provided by the World Bank, three separate imputed panel datasets were created to train their model. Study revealed a positive correlation between the adoption of digital payment platforms and GDP economic growth in the 189 nations sampled from the World Bank dataset, a 1\% increase in digital payment adoption boosted GDP growth on average by 6 - 8\% [5].

        Edquist, Goodridge, and Haskel investigated and detrmined that the economic value generated by streaming services and argues that much of this value is not fully captured by conventional GDP measures. The authors relied on economic measurement frameworks and conceptual analysis rather than algorithmic or machine learning models. Through their work, identified a gap in existing national accounting systems’ ability to quantify digital service value, suggesting the need for improved measurement methods and complementary analytical approaches [6].

        Magoutas, Chaideftou, Skandali, and Chountalas report that the Information and Communication Technology (ICT) development, measured through indicators such as: digital infrastructure, human capital, and technology integration—affects GDP growth across European Union countries [7]. The authors use structural equation modeling to assess the relationship between ICT advancement and economic performance. Robust data was sourced from official, licensed databases such as the Digital Economy and Society Index (DESI), Statistical Office of the European Union (EUROSTAT), and the Organization for Economic Cooperation and Development (OECD). SmartPLS version 4.0.9.2 was chosen as the software to perform analysis, providing estimation of ordinary least squares and reliable examination of influences and interdependencies in the datasets. Limitations are related to regional scope and the use of traditional statistical models, highlighting the need for predictive and nonlinear approaches that can better capture the complex relationships of digital-driven growth.
        
        A review of 64 peer-reviewed journal articles -- selected from an initial pool of 2,977 studies using PRISMA guidelines and a hybrid quantitative-qualitative methodology revealed that 75\% of empirical research used multivariate models that incorporated forward-looking indicators such as search engine trends and economic sentiment measures[8]. LSTM, Random Forest, and Gradient Boosting methods have been identified as the most effective algorithms, 40\% of studies reporting ANN-based models as the best performing [9]. Key implementation strategies include demand decomposition into baseline and volatile components, rolling window validation techniques, and volatility-based model selection using the Coefficient of Variation. The primary limitations reported were: data quality concerns, limited generalizability, interpretability challenges introduced from black-box models, and the disconnect from theoretical framework to real-world volatility.

        Findings from 21 studies examining ML applications for economic forecasting through the lens of Sustainable Development Goal 8 (Decent Work and Economic Growth) evaluates methodologies across three dimensions:  economic (GDP, inflation, exchange rates), social (unemployment, employability), and environmental (carbon dioxide emissions, climate variables, renewable energy adoption) [10]. Predominant techniques include Gradient Boosting (XGBoost, LightGBM), Random Forest, Support Vector Machines, and ARIMA, with RMSE, MAE, and MAPE serving as the primary evaluation metrics. A significant research gap emerges from the lack of integrated modeling frameworks that simultaneously address economic growth and decent work outcome. Existing studies typically analyze macroeconomic performance or labor market indicators in isolation, overlooking the dynamic interdependencies between these domains.
        
        Guarino, Montieri, Ciunzo, and Pescape find that their model reported an underestimation of the energy sector and mixed results in other sectors. Prediction errors also varied, having a higher error rate for companies with more than 50 employees [11].Applying four supervised ML models, Linear Regression, Decision Tree, Random Forest Regressor, and XGBoost the study input the following indicators in order to make their predictions: sales revenue, production cost, EBITDA (earnings before interest, taxes, depreciation, and amortization), net profit, total assets, shareholder's equity, and net financial position. Their dataset comprised of 529 companies with 53 distinct features, it should also be noted that the geographical region within Italy that the company was located also affected the prediction outcome, providing the majority of errors with southern based companies and consistent results with northern and central based companies.

\newpage
\bibliographystyle{plain}
\bibliography{references}

\begin{thebibliography}{00}
\bibitem{b1} M.D. Adewale, D.U. Ebem, O. Awdodele, A. Sambo-Magaji, E.M. Aggrey, E.A. Okechalu, R.E. Donatus, K.A. Olayanju, A.F. Owolabi, J.U. Oju, O.C. Ubadike, G.A. Otu, U.I. Muhammed, O.R. Danjuma, O.P. Oluyide, ''Predicting gross domestic product using the ensemble learning method'', Systems and Soft Computing 6, 2024.
\bibitem{b2} A. Richardson, T. van Florenstein Mulder, ''Nowcasting GDP using machine-learning algorithms: A real-time assessment,'' in International Journal of Forecasting, 2021, pp. 941--48.
\bibitem{b3} I. Vasenska, ''Comparative analysis of machine learning and deep learning models for tourism demand forecasting with economic indicators'', in FinTech, volume 4, Issue 3, 2025.
\bibitem{b4} L. Vanscura, T. Tatay, T. Bareith, ``Navigating AI-driven financial forecasting: a systematic review of current status and critical research gaps,'' in Forecasting, vol. VII, Issue 3, 2025.
\bibitem{b5} A. Birigozzi, C. De Silva, P. Luitel, ''Digital payments and GDP growth: A behavioural quantitative analysis,'', in Research in International Business and Finance, vol. 75, 2025.
\bibitem{b6} H. Edquist, P. Goodridge, J. Haskel, ''The economic impact of streaming beyond GDP'', Applied Economics Letters, vol. 29, Issue 5, pp. 403 - 408, 2022.
\bibitem{b7} A. Magoutas, M. Chaideftou, D. Skandali, P. Chountalas, ''Digital progression and economic growth: analyzing the impact of ICT advancements on the GDP of European Union countries,'' in Economies, vol. XII, Issue 3, 2024.
\bibitem{b8} M. Muth, M. Lingenfelder, G. Nufer, ``The application of machine learning for demand predicting under macroeconomic volatility a systematic literature review,'' in Management Review Quarterly, vol. 75, pp. 2759 - 2802, 2025.
\bibitem{b9} I. Elbatal, M. Sarwar, F. Jamal, M. Daniyal, Z. Hussain, A. Ben Ghorbal, ''Modelling on gross domesting product annual growth rate data by using time series, machine learning, and probability models'', in Journal of Radiation Research and Applied Sciences, vol. 18, Issue 2, 2025.
\bibitem{b10} M. Tamer Shaekr, H. El-Batatnouny, M. Abdelsalam, Y. Helmy, ''Systematic review of machine learning approaches in forecasting economic growth and enhancing decent work opportunities: a comprehensive analysis'', 2024 6th International Conference on Computing and Informatics (ICCI), New Cairo - Cairo, Egypt, 2024, pp. 377-392.
\bibitem{b11} I. Guarino, A. Montieri, D. Ciunozo, A. Pescape, ''Explainable predictive analysis of the economic and financial performance of Italian publicly owned companies'', 2025 International Conference on Artificial intelligence in Information and Communication (ICAIIC), Fukuoka, Japan, 2025, pp. 0087 - 0092.
\end{thebibliography}
\end{document}